\documentclass[a4paper]{article} %

\date{} %

% Please, do not change any of the above lines

\usepackage{amssymb} % add other LaTeX packages you need here.

\begin{document}
\setcounter{page}{1} % Please, do not delete this line

% Your title
\title{Dealing with {\rm C}$^\star$-algebras by means of \\ algebra and order}

% Authors
\author{Thomas Vetterlein \\ %
       Johannes Kepler University Linz \\ % Affiliation 1
       \texttt{Thomas.Vetterlein@jku.at} % Only one corresponding e-mail, please
       }%

\maketitle

Considerable effort has been devoted to reconstructing complex Hilbert spa\-ces -- the foundational models of quantum physics -- from simpler structures such as ortholattices or orthosets. It was shown, for instance, that in a sense, a Hilbert space can be reduced to its orthogonality relation. A similar challenge arises for a class of mathematical models actually much better tailored to the needs of quantum mechanics: {\rm C}$^\star$-algebras. While a {\rm C}$^\star$-algebra likewise gives rise to a notion of orthogonality, this relation alone does not provide much of information, unlike in the Hilbert space case. It remains unclear how to fully characterize a {\rm C}$^\star$-algebra through purely algebraic or order-theoretic means.

The unsatisfactory situation has motivated our attempt to develop the basic theory of {\rm C}$^\star$-algebras with a primary focus on algebraic and order-theoretic aspects. While standard approaches employ methods of analysis and set-theoretic topology from the outset, a considerable portion of the theory can actually be developed without spectral theory. The order unit space of self-adjoint elements, the unitary group, approximate units, faces, and ideals can all be treated using methods limited to algebraic reasoning and the convergence of power series. 

On this basis, we may then develop Gelfand theory following the philosophy of pointless topology: rather than a locally compact Hausdorff space, the space associated with a commutative {\rm C}$^\star$-algebra is identified as a continuous regular frame. Ultimately, this approach allows us to represent elements of a {\rm C}$^\star$-algebra as spectral maps, in analogy to the case of operators on a Hilbert space.

Constructive Gelfand duality for {\rm C}$^\star$-algebras was developed by Banaschewski and Mulvey \cite{BaMu} as well as by Coquand and Spitters \cite{CoSp}, with the non-unital case subsequently studied by Henry \cite{Hen}. While our approach is possibly not constructive, it similarly avoids the (full) Axiom of Choice. Moreover, the procedure applies not only to an isolated commutative {\rm C}$^\star$-algebra but also to a {\rm C}$^\star$-subalgebra embedded within a larger {\rm C}$^\star$-algebra. Consequently, the structural relationships between the constituents of commutative {\rm C}$^\star$-subalgebras and those of the ambient algebra become to some extent transparent.

{\bf Acknowledgement.} This research was funded in part by the Austrian Science Fund (FWF) 10.55776/ PIN5424624 and the Czech Science Foundation (GA\v CR) 25-20013L.

\begin{thebibliography}{BaMu2}

\bibitem{BaMu} B. Banaschewski, C. Mulvey,
A globalisation of the Gelfand duality theorem,
{\it Ann. Pure Appl. Logic} {\bf 137} (2006), 62--103.

\bibitem{CoSp} T. Coquand, B. Spitters,
Constructive Gelfand duality for {\rm C}$^\star$-algebras,
{\it Math. Proc. Camb. Philos. Soc.} {\bf 147} (2009), 339--344.

\bibitem{Hen} S. Henry,
Constructive Gelfand duality for non-unital commutative C*-algebras,
arXiv:1412.2009 (2014).

\end{thebibliography}

\end{document}
